\studynote{2}{Tue 03 Dec 2024 17:09}{SLLN for BRWRE}

\subsubsection{Rigorous argument for BRW-RE}

In this section we prove lower bound for Proposition~\ref{prop:SLLN-BRWRE}. The proof takes several steps:

\textbf{Step 1. Construction of coarse-grained tree.}
%Start from a vertex $v \in D_m$ where $S_v = x$, we define a tree $T^v$ as a subtree of the original tree structure. Take $T^v_0 = \{v\} $. Let $D^T_w$ denote the descendants of $w$ in the tree $T$. Then
%\[
%	D^{T}_w := \{w' \in D_w: \rho(w', w) = L, S_{w'} > S_w + (x^* - \varepsilon) L \} 
%.\] 

%Alternatively we can view $T$ as a percolated subtree of the coarse-grained tree $T_C$, with the edge condition $S_{w'} > S_w + (x^* - \varepsilon) L$.

Let $T$ be a tree. We denote by $T(v)$ the subtree rooted at vertex $v$. We denote by $T(k)$ the set of vertexes in the $k$-th generation. We denote $T(v, k) = T(v)(k)$.

Recall $\mathcal{T}$ is the $k$-ary tree associated with the BRW starting at $0$. Start from any $k$-ary subtree $T$ of $\mathcal{T}$ we construct $L$\textit{-coarse-grain} of $T$, denoted as $C_L(T)$, as follows:
\[
	C_L(T)(0) = T(0), \quad C_L(T)(v, 1) = T(v, L)
.\] 

On the coarse-grained tree, we define a percolation condition: for $w \in C_L(T)(v, 1)$,
\[
	w \leftrightarrow v \iff  S_w \ge S_v + x L
.\] 
where $x = x^* - \varepsilon$.

Denote by $\tilde T$ the percolated subtree of $C_L(T)$.



\textbf{Step 2. Realign the tree via coupling}. 
%Take a tree $T_C$. Now, for every element $v \in T_C(1)$, consider a $k$-ary tree $K_v$ isomorphic to the subtree $T_C(v)$. Assign edge values to $K_v$ iteratively as follows: if $S^{K_v}_w < S^{T_C}_w$ then 

Consider a $k$-ary subtree $T$ of $\mathcal{T}$. We define the $L$-\textit{alignment} of $T$, denoted as $A_L(T)$, as follows:

$A_L(T)(0), \ldots, A_L(T)(L - 1)$ is chosen to be identical to $T(0), \ldots, T(L - 1)$. For each $v \in T(L)$, consider an independent copy of BRW starting from $2 \left\lfloor x L / 2 \right\rfloor$, giving a $k$-ary tree $K_v$. Pick a tree isomorphism $\varphi: K_v \to  T(v)$.

Now, take $A_L(T)(L) = \{ K_v(0): v \in T(L)\} $, with $K_v(0)$ being the descendent of the ancestor of $v$. For each $w \in A_L(T)(L)$, if $S_w < S_{\varphi(w)}$, take $A_L(T)(w, 1) = K_v(w, 1)$ and continue the construction on each element of $A_L(T)(w, 1)$. Otherwise, take $A_L(T)(w) = T(w)$.

The resulting construction is a coupling between the original BRW, and a new process with all the elements shifted to position  $xL$ by time $L$, and if the original BRW is on the right of $xL$, the new BRW never exceeds original BRW. That is, $A_L(T)$ represents a BRW that goes strictly slower than $T$.

Now iterate the construction, and define the $L$-\textit{superalignment} of $T$, denoted as $A^*_L(T)$, as
\[
	A^*_L(T) = \lim_{n \to \infty } A_{n L} \ldots A_{L} (T)
.\] 

From now on, we focus on the tree $\widetilde{A^*_L}(\mathcal{T}(v))$, which we simply denote by $\tilde T$ for short.

\textbf{Step 3. Super-criticality of $\tilde T$.} We apply the following theorem to show that $T^v$ is supercritical for any $v$ and large enough $L$:

\begin{theorem}[Athreya-Karlin Theorem '1970]
	\label{thm:AthreyaKarlthm}
	Let $\xi_0, \xi_1, \ldots$ be a stationary ergodic process of probability generating functions for offspring distribution of a branching process. Then the branching process survives with positive probability if the following two conditions hold:
	\begin{enumerate}
		\item $E \left[ - \log \left( 1 - \varphi_{\xi_0}(0) \right)  \right] < \infty $,
		\item $E \left[ \log \phi'_{\xi_0}(1) \right] ^- < E \left[ \log \varphi'_{\xi_0}(1) \right] ^+ \le \infty $.
	\end{enumerate}
\end{theorem}

Condition (1) holds automatically due to uniform boundedness of $\omega$. To verify (2),
It suffices to show that
\[
	\mathbb{E}\left[ \log E_\omega\left[ |\tilde T(1)| \right]  \right] > 0, \quad \text{ for large }L
,\] 
which is equivalent to
\[
	\mathbb{E}\left[ \log P_\omega(X_L > xL) \right] > - L \log k, \quad \text{ for large }L
.\] 

To show this, we use the following LDP:
\begin{theorem}[Position LDP for ERWRE, Comets et al., '2000]
	\label{thm:PositionLDPERWREComets}
	(Theorem 1.) Assume $\eta \in M_1^e(\Sigma)^K$. For $\eta$-a.e. $\omega$, the distributions of $X_n / n$ under $P_\omega$ satisfy a large deviation principle with convex, good rate function $I_\eta^q$.
\end{theorem}

We have, from large deviation principle,
\[
	\frac{1}{L} \log P_\omega(X_L > x L) = - I(x) + o_{L, \omega}(1)
.\] 

Uniform ellipticity implies that the error term $o_{L, \omega}$ is deterministically bounded. Applying DCT yields
\[
	\mathbb{E}\left[\frac{1}{L} \log P_\omega(X_L > x L)\right] = - I(x) + o_{L}(1)
.\]

Since $I(x) < \log k$,
\[
	\mathbb{E}\left[\log P_\omega(X_L > x L)\right] > - L \log k \quad \text{ for large }L
.\] 

\textbf{Step 4. Many many trees.}
%We pick the particle $v_1$ that hits $1$ first to start a tree $T^{v_1}$. We then pick a particle $v_2$ that hits $2$ first, and which is not a descendent of $v_1$, to start a tree $T^{v_2}$. Continue in this fashion, we obtain a countable sequence of independent coarse-grained trees. By Borel-Cantelli, there must exist a tree $T^w$ that survives. Then
%\[
%	\liminf S_{k L}/ k L \ge (x^* - \varepsilon)
%.\] 

Pick $v_1 \in \mathcal{T}(1)$, $v_2 \in \mathcal{T}(2) \setminus \mathcal{T}(v_1)$, $v_3 \in \mathcal{T}(3) \setminus \left( \mathcal{T}(v_1) \cup \mathcal{T}(v_2) \right) $, etc. Then $\mathcal{T}(v_1), \mathcal{T}(v_2), \ldots$ are independent. Therefore, $\widetilde{A^*_L}(\mathcal{T}(v_i))_{i= 1, 2, \ldots}$ are independent, and each have a positive probability, independent of $i$, to survive.

By Borel-Cantelli, there exists $i$ such that $\widetilde{A^*_L}(\mathcal{T}(v_i))$ survives. Take a surviving path $w_k \in \widetilde{A^*_L}(\mathcal{T}(v_i))(k)$. Then
\[
	M_{i + kL} \ge S_{w_k} = S_{v_i} + k (xL)
.\] 

Hence
\[
	\liminf_{n \to \infty } \frac{M_n}{n} \ge x
.\] 

This concludes the proof of lower bound for BRWRE-SLLN.
